\documentclass[main]{subfiles}
\usepackage[sols]{workbook}
\numbering{ws1-3}

\begin{document}

\ifSubfilesClassLoaded{%
  \chaptermark{\chapterone}
}{}

\section{Area and Volume} \label{ws1-3}


\begin{framed}
\textbf{Lesson Goals:}
\begin{itemize}
\item You should be able to sketch regions in the $xy$-plane bounded by basic curves, and find the bounds of the region.
\item You should be able to use our general slicing strategy (\textit{Slice} $\rightarrow$ \textit{Approximate} $\rightarrow$ \textit{Sum} $\rightarrow$ \textit{Limit}) to write and evaluate a definite integral that gives the area bounded by basic curves in the $xy$-plane.
\item You should be able to decide whether to slice a region vertically or horizontally to find its area most easily and efficiently, based on the geometry of the region and the resulting integral(s).
\item You should be able to evaluate challenging definite integrals by interpreting them as areas and finding them via slicing.
\item You should be able to extend our general slicing strategy for computing volumes when the volume of each slice is simple.
\end{itemize}
\end{framed}

\begin{enumerate}

\item \label{area0}
  \begin{problem}[\vfill]
        Let $\mathcal{R}$ be the region enclosed by the curves $y = x^2 - 1$
  and $y = x + 1$. Find the area of this region using our slicing strategy.
  \end{problem}


    \begin{solution}
      To find the $x$-coordinates of the points of intersection, we solve
      $x^2 - 1 = x + 1$:
      \begin{align*}
        x^2 - 1 &= x + 1 \\
        x^2 - x - 2 & = 0 \\
        (x - 2)(x + 1) & = 0 \\
        x &= -1, 2
      \end{align*}
      So, the points of intersection are \fbox{$(-1, 0)$ and $(2, 3)$}.
      (To find the $y$-coordinate, we plug the $x$-coordinate into $y =
      x^2 - 1$ or $y = x + 1$.) 
 
      The region has width 3 (from $x = -1$ to $x = 2$), and we're
      dividing it into $n$ pieces, so $\Delta x = \frac{3}{n}$.

      $x_k$ is $k$ ``slice widths'' to the right of $x_0$, so $x_k = x_0 +
      k \Delta x$.  We know that $x_0 = -1$, so $\boxed{x_k = -1 + k
        \Delta x}$.  (Notice that then $x_n = -1 + n \Delta x = -1 + n
      \cdot \frac{3}{n} = 2$, which matches our picture.)
  


      We are approximating the area of the $k$-th slice by pretending that
      the slice is really the rectangle shown in the picture below.  The width
      of the rectangle is the same as the width of the slice, which is
      $\Delta x$.

      The rectangle shares the same right side as the slice, so let's
      figure out the coordinates of the top and bottom points of the
      right side of the slice (the blue points in the following picture).
      \begin{center}
        \includegraphics[width=1.5in]{ws/images/area0-slice-approximate-points.pdf}
      \end{center}      
      We know that the right end of the slice
      has $x = x_k$.  The bottom of the slice is the curve $y = x^2 -
      1$, so the $y$-coordinate of the bottom blue point is $x_k^2 - 1$.
      That is, the bottom blue point has coordinates $(x_k, x_k^2 - 1)$.
      Similarly, the top blue point has coordinates $(x_k, x_k + 1)$.  The
      height of the rectangle is simply the $y$-coordinate of the top blue
      point minus the $y$-coordinate of the bottom blue point, or $(x_k +
      1) - (x_k^2 - 1)$.

      So, the area of the rectangle is $\boxed{[(x_k + 1) - (x_k^2 - 1)]
        \Delta x}$. 
  
      To get the approximate area of the whole region, we sum the
      approximate areas of all of the slices.  Therefore, the approximate
      area of the region is
      \[ \boxed{[(x_1 + 1) - (x_1^2 - 1)] \Delta x + [(x_2 + 1) - (x_2^2 -
        1)] \Delta x + \cdots + [(x_n + 1) - (x_n^2 - 1)] \Delta x}\] This
      can also be written in summation notation as $\displaystyle
      \boxed{\sum_{k=1}^n [(x_k + 1) - (x_k^2 - 1)] \Delta x}$.
    
        We can take the appropriate limit of the previous sum and the resulting integral is $$\displaystyle\int_{-1}^2[(x+1)-(x^2-1)]\,dx$$
    \end{solution}

\pspace{\newpage}
\item  \label{area-ln-horizontal}

  
  \begin{problem}[\vfill]
    Find the area under the curve $y = \ln x$ from $x = 1$ to $x = e$.
  \end{problem}

  \begin{solution}
    Here is the area we are interested in:
    \begin{center}
      \includegraphics[height=1in]{ws/images/area-ln-horizontal}
    \end{center}

    We know that this area is $\displaystyle \int_1^e \ln x~dx$, but that
    integral is a little hard to compute (you may not know how to do it
    yet, which is fine --- we'll talk about it next week).  Let's try
    another strategy.

    Let's \textbf{slice} the region into \textit{horizontal} strips of
    equal height $\Delta y$ (so we slice the interval $[0, 1]$ on the
    $y$-axis), label $y_0 = 0, \ldots, y_n = 1$ as usual, and focus on the
    $k$-th slice.
    \begin{center}
      \includegraphics[height=1in]{ws/images/area-ln-horizontal-slice}
    \end{center}
    The curve is $y = \ln x$, which we can also think of as $x = e^y$.
    We can \textbf{approximate} the slice by a rectangle (shown in
    purple):    
    \begin{center}
      \includegraphics[height=1in]{ws/images/area-ln-horizontal-approximating-rectangle}
    \end{center}
    Since the curve is $x = e^y$, the top left point of the rectangle has
    coordinates $\left( e^{y_k}, y_k \right)$.  The top right point has
    coordinates $(e, y_k)$.  Therefore, the width of the rectangle is $e -
    e^{y_k}$.  The rectangle has height $\Delta y$, so its area is $(e -
    e^{y_k}) \Delta y$.  So, we've shown that the area of the $k$-th slice
    is approximately $(e - e^{y_k}) \Delta y$.  \textbf{Summing} over all
    slices and \textbf{taking the limit} as the number of slices tends to
    infinity, we find that the area is $\displaystyle \lim_{n \to \infty}
    \sum_{k=1}^n (e - e^{y_k}) \Delta y$, which is equal to the integral
    $\displaystyle \int_0^1 (e - e^y)~dy$.  We can easily evaluate this
    integral; it is equal to $\displaystyle ey - e^y \Big|_0^1 = (e - e) -
    (0 - 1) = \boxed{1}$.
  \end{solution}

%\probonly{\newpage}

\item  \label{area-3curves-3}

  \note{If students get this quickly, encourage them to also think about
    what happens if they slice ``the other way'' (they should be able to
    slice either horizontally or vertically).  Also, if students just
    write down an integral directly, ask them how they sliced and what's
    happening with the $k$-th slice.}
  
  \begin{problem}[\vfill]
    Find the area enclosed by $y = -2x$, $y = x^3$, and $y = 1$.
  \end{problem}

  \begin{solution}
    Let's draw a picture so we can see what area we're talking about.  If
    we graph $y = -2x$, $y = x^3$, and $y = 1$, we see that the area we're
    talking about is this:
    \begin{center}
      \includegraphics[width=1.5in]{ws/images/area-3curves-3-region}
    \end{center}
    As always, the first step is to \textbf{slice}.  We can either slice
    horizontally or vertically: 
    \begin{center}
      \includegraphics[width=1.5in]{ws/images/area-3curves-3-hslice}
      \hspace{1in}
      \includegraphics[width=1.5in]{ws/images/area-3curves-3-vslice}
    \end{center}
    Notice that, if we use vertical slices, there are two ``types'' of
    slices --- slices to the left of the $x$-axis go between the curves $y
    = -2x$ and $y = 1$ while slices to the right of the $x$-axis go
    between $y = x^3$ and $y = 1$.  This means we need to consider these
    two areas separately.
    
    If we use horizontal slices, then all of the slices have the same
    description: the left side is on $y = -2x$ (or $x = -y/2$), and the
    right side is on $y = x^3$ (or $x = y^{1/3}$).  So, using horizontal
    slices is slightly more efficient.  If we slice $[0, 1]$ into $n$
    pieces and label $y_0 = 0, \ldots, y_n = 1$ as usual, then we can
    \textbf{approximate} the area of the $k$-th slice as $\left(y_k^{1/3}
      + \frac{y_k}{2} \right) \Delta y$.  \textbf{Summing} and
    \textbf{taking the limit} as $n \to \infty$ gives us $\displaystyle
    \lim_{n \to \infty} \sum_{k=1}^n \left( y_k^{1/3} + \frac{y_k}{2}
    \right) \Delta y$, or $\displaystyle \int_0^1 \left( y^{1/3} +
      \frac{y}{2} \right)~dy = \left. \frac{3}{4} y^{4/3} + \frac{1}{4}
      y^2 \right|_0^1 = \boxed{1}$.

    \textit{Note:} If you use vertical slices, the correct sum of
    integrals is $\displaystyle \int_{-1/2}^0 (1 + 2x)~dx + \int_0^1 (1 -
    x^3)~dx$.
  \end{solution}
  
  %\probonly{\newpage}

\pspace{\newpage}


\item \label{area-arcsin-arccos}
 

  True or false.  Explain your reasoning thoroughly in each.
  
  \begin{enumerate}

  \item
    \begin{problem}[\vfill]
      $\displaystyle \int_0^1 \arccos x~dx = \int_0^{\pi / 2} \cos y~dy$.
    \end{problem}

    \begin{solution}
      This is \fbox{true}.  One way to see this is simply to evaluate the
      two integrals.

      Alternatively, we can interpret $\displaystyle \int_0^1 \arccos
      x~dx$ as the area under $y = \arccos x$ from $x = 0$ to $x = 1$,
      which is the area of the following region:
      \begin{center}
        \includegraphics[width=1.5in]{ws/images/area-arcsin-arccos-arccos}
      \end{center}
      Suppose we instead decided to find the area by slicing horizontally
      (as you did in Problem Set 3, {\#3}).  Then, we'd slice the interval
      $\left[0, \frac{\pi}{2} \right]$ (on the $y$-axis) into $n$ pieces
      of equal thickness $\Delta y$ and label $y_0 = 0, \ldots, y_n =
      \frac{\pi}{2}$.
      \begin{center}
        \includegraphics[width=1.5in]{ws/images/area-arcsin-arccos-arccos-hslices}
      \end{center}
      Now, the $k$-th slice has height $\Delta y$ and approximate width
      $\cos y_k$ (since the curve is $y = \arccos x$, or $x = \cos y$), so
      the area of the $k$-th slice is approximately $\cos y_k \Delta y$.
      Summing over all slices and taking the limit as $n \to \infty$, the
      area is indeed given by $\displaystyle \int_0^{\pi / 2} \cos y~dy$. 
    \end{solution}
  
  \item \note{This seems to be a somewhat common mistake that happens when
      students don't draw a picture.}

    \begin{problem}[\vfill]
      $\displaystyle \int_0^1 \arcsin x~dx = \int_0^{\pi / 2} \sin y~dy$.
    \end{problem}

    \begin{solution}
      This is \fbox{false}.  We can use the same reasoning as in the
      previous part to see that $\displaystyle \int_0^1 \arcsin x~dx$ is
      the blue region in the following picture (the area under $y =
      \arcsin x$ from $x = 0$ to $x = 1$), while $\displaystyle
      \int_0^{\pi / 2} \sin y~dy$ is the area of the purple region:
      \begin{center}
        \includegraphics[width=1.5in]{ws/images/area-arcsin-arccos-arcsin}
      \end{center}
      The two regions clearly have different areas, so the statement is
      false. 
    \end{solution}

  \end{enumerate}
\pspace{\newpage}
\item

\note{Just so that they have one problem about volumes which is not about solids of revolution, similarly to what they have on the PSet.}

\begin{problem}
    Use slicing to obtain the volume of a square pyramid with height 5 and base with sides of length 2. You can verify your answer is correct by using the formula for the volume of a pyramid.
\end{problem}
\probonly{\vfill}
\begin{solution}
    We will call $x$ the vertical distance measured perpendicularly from the base of the pyramid, as indicated in the picture below:

\begin{center}
      \includegraphics[width=1.5in]{ws/images/Solid_Pyramid_noslice.png}
\end{center}

We will \textbf{slice} the pyramid parallel to the base, since we know that shape is a square. Let us define $\Delta x = 5/n$ and $x_k = k \Delta x$, so that $x_0 = 0, x_n = 5$. A generic $k$-th slice is indicated in the picture below:

\begin{center}
      \includegraphics[width=1.5in]{ws/images/Solid_Pyramid_slice.png}
\end{center}

We now \textbf{approximate} the volume of the slice as the volume of a box with height $\Delta x$ and a square base with side length $l_k$. That means the volume of the slice is approximately $(l_k)^2 \Delta x$. To find $l_k$ we can use proportions on the side triangle of the pyramid, as indicated by the figure

\begin{center}
      \includegraphics[width=1.5in]{ws/images/Solid_Pyramid_slice_front.png}
\end{center}

We obtain the following relation:
\[\frac{5-x_k}{l_k} = \frac{5}{2}\]
It follows that $l_k = \frac{2}{5}(5-x_k)$ and the volume of the $k$-th slice is approximately $\frac{4}{25}(5-x_k)^2 \Delta x$. 

\textbf{Summing} over all slices we get the approximate volume for the pyramid
$\displaystyle \sum_{k=1}^n\frac{4}{25}(5-x_k)^2 \Delta x$ and \textbf{taking the limit} as the number of slices tend to infinity we get the exact volume of the pyramid, given by the integral
\[\int_0^5 \frac{4}{25}(5-x)^2 \,dx \]
We can solve this integral by substitution, making $u = 5-x$, so that 
\[\int_0^5 \frac{4}{25}(5-x)^2 \,dx = -\int_{5}^0 \frac{4}{25}u^2 \,du = \frac{4}{3\cdot 25} \left.u^3\right|_0^5 = \frac{4\times 5}{3} \]
Which matches the result of the formula for the volume of the pyramid as $\frac{(\text{Area of the base})(\text{Height})}{3}$.
    
\end{solution}

\end{enumerate}
\end{document}